\section[邶]{邶}

\subsection[柏舟]{柏舟}

汎彼柏舟，亦汎其流。耿耿不寐，如有隱憂。微我無酒，以敖以遊。

我心匪鑒，不可以茹。亦有兄弟，不可以據。薄言往愬，逢彼之怒。

我心匪石，不可轉也。我心匪席，不可卷也。威儀棣棣，不可選也。

憂心悄悄，愠于群小。覯閔既多，受侮不少。静言思之，寤辟有摽。

日居月諸，胡迭而微？心之憂矣，如匪澣衣。静言思之，不能奮飛！

\subsection[緑衣]{緑衣}

緑兮衣兮，緑衣黄裏，心之憂矣，曷維其已。

緑兮衣兮，緑衣黄裳，心之憂矣，曷維其亡。

緑兮絲兮，女所治兮，我思古人，俾無訧兮。

絺兮綌兮，淒其以風，我思古人，實獲我心。

\subsection[燕燕]{燕燕}

燕燕于飛，差池其羽，之子于歸，遠送于野，瞻望弗及，泣涕如雨。

燕燕于飛，頡之頏之，之子于歸，遠于將之，瞻望弗及，佇立以泣。

燕燕于飛，下上其音，之子于歸，遠送于南，瞻望弗及，實勞我心。

仲氏任只，其心塞淵，終温且惠，淑慎其身，先君之思，以勗寡人。

\subsection[日月]{日月}

日居月諸，照臨下土，乃如之人兮，逝不古處，胡能有定，寧不我顧。

日居月諸，下土是冒，乃如之人兮，逝不相好，胡能有定，寧不我報。

日居月諸，出自東方，乃如之人兮，德音無良，胡能有定，俾也可忘。

日居月諸，東方自出，父兮母兮，畜我不卒，胡能有定，報我不述。

\subsection[終風]{終風}

終風且暴，顧我則笑，謔浪笑敖，中心是悼。

終風且霾，惠然肯來，莫往莫來，悠悠我思。

終風且曀，不日有曀，寤言不寐，願言則嚏。

曀曀其陰，虺虺其靁，寤言不寐，願言則懷。

\subsection[擊鼓]{擊鼓}

擊鼓其鏜，踴躍用兵，土國城漕，我獨南行。

從孫子仲，平陳與宋，不我以歸，憂心有忡。

爰居爰處，爰喪其馬，于以求之，于林之下。

死生契闊，與子成説，執子之手，與子偕老。

于嗟闊兮，不我活兮，于嗟洵兮，不我信兮。

\subsection[凱風]{凱風}

凱風自南，吹彼棘心；棘心夭夭，母氏劬勞。

凱風自南，吹彼棘薪；母氏聖善，我無令人。

爰有寒泉，在浚之下；有子七人，母氏勞苦。

睍睆黄鳥，載好其音；有子七人，莫慰母心。

\subsection[雄雉]{雄雉}

雄雉于飛，泄泄其羽，我之懷矣，自詒伊阻。

雄雉于飛，下上其音，展矣君子，實勞我心。

瞻彼日月，悠悠我思，道之云遠，曷云能來。

百爾君子，不知德行，不忮不求，何用不臧。

\subsection[匏有苦葉]{匏有苦葉}

匏有苦葉，濟有深涉，深則厲，淺則揭。

有瀰濟盈，有鷕雉鳴，濟盈不濡軌，雉鳴求其牡。

雝雝鳴鴈，旭日始旦，士如歸妻，迨冰未泮。

招招舟子，人涉卬否，人涉卬否，卬須我友。

\subsection[谷風]{谷風}

習習谷風，以陰以雨。黽勉同心，不宜有怒。

采葑采菲，無以下體？德音莫違，及爾同死。

行道遲遲，中心有違。不遠伊邇，薄送我畿。

誰謂荼苦，其甘如薺。宴爾新昏，如兄如弟。

涇以渭濁，湜湜其沚。宴爾新昏，不我屑以。

毋逝我梁，毋發我笱。我躬不閲，遑恤我後。

就其深矣，方之舟之。就其淺矣，泳之游之。

何有何亡，黽勉求之。凡民有喪，匍匐救之。

不我能慉，反以我爲讎。既阻我德，賈用不售。

昔育恐育鞫，及爾顛覆。既生既育，比予于毒。

我有旨蓄，亦以御冬。宴爾新昏，以我御窮。

有洸有潰，既詒我肄。不念昔者，伊余來塈。

\subsection[式微]{式微}

式微式微，胡不歸？微君之故，胡爲乎中露？

式微式微，胡不歸？微君之躬，胡爲乎泥中？

\subsection[旄丘]{旄丘}

旄丘之葛兮，何誕之節兮，叔兮伯兮，何多日也。

何其處也，必有與也，何其久也，必有以也。

狐裘蒙戎，匪車不東，叔兮伯兮，靡所與同。

瑣兮尾兮，流離之子，叔兮伯兮，褎如充耳。

\subsection[簡兮]{簡兮}

簡兮簡兮，方將萬舞，日之方中，在前上處，碩人俣俣，公庭萬舞。

有力如虎，執轡如組，左手執籥，右手秉翟，赫如渥赭，公言錫爵。

山有榛，隰有苓，云誰之思，西方美人，彼美人兮，西方之人兮。

\subsection[泉水]{泉水}

毖彼泉水，亦流于淇，有懷于衞，靡日不思，孌彼諸姬，聊與之謀。

出宿于泲，飲餞于禰，女子有行，遠父母兄弟，問我諸姑，遂及伯姊。

出宿于干，飲餞于言，載脂載舝，還車言邁，遄臻于衞，不瑕有害。

我思肥泉，兹之永歎，思須與漕，我心悠悠，駕言出遊，以寫我憂。

\subsection[北門]{北門}

出自北門，憂心殷殷，終窶且貧，莫知我艱，已焉哉，天實爲之，謂之何哉。

王事適我，政事一埤益我，我入自外，室人交徧讁我，已焉哉，天實爲之，謂之何哉。

王事敦我，政事一埤遺我，我入自外，室人交徧摧我，已焉哉，天實爲之，謂之何哉。

\subsection[北風]{北風}

北風其涼，雨雪其雱，惠而好我，攜手同行，其虚其邪，既亟只且。

北風其喈，雨雪其霏，惠而好我，攜手同歸，其虚其邪，既亟只且。

莫赤匪狐，莫黑匪烏，惠而好我，攜手同車，其虚其邪，既亟只且。

\subsection[静女]{静女}

静女其姝，俟我於城隅，愛而不見，搔首踟躕。

静女其孌，貽我彤管，彤管有煒，説懌女美。

自牧歸荑，洵美且異，匪女之爲美，美人之貽。

\subsection[新臺]{新臺}

新臺有泚，河水瀰瀰，燕婉之求，蘧篨不鮮。

新臺有洒，河水浼浼，燕婉之求，蘧篨不殄。

魚網之設，鴻則離之，燕婉之求，得此戚施。

\subsection[二子乘舟]{二子乘舟}

二子乘舟，汎汎其景，願言思子，中心養養。

二子乘舟，汎汎其逝，願言思子，不瑕有害。
